\chapter{Conclusion}
\todome{DRAFT START (Claude 04.08): conclusion chapter, five paragraphs --- restate, diagnosis, NGT, ASL Citizen, two-regime takeaway + outlook; no numbers beyond those already in the body}
This thesis asked whether sign language recognition models lean on the
incidental appearance of the recorded signers (clothing, skin tone,
body shape, environment) rather than the signing itself, and whether
synthetic appearance counterfactuals, edited copies of training videos in
which appearance changes while the performed sign is preserved, can
remove that bias. We pursued the question through one diagnosis and two
interventions: motion-locked avatar re-renders of a
self-recorded NGT744 corpus, where the counterfactual property holds by
construction, and generative per-frame appearance edits of ASL
Citizen videos, the scalable approximation of the same idea at benchmark
scale.
% Restatement condenses the abstract (lines 76-88) and the introduction's
% pipeline framing (lines 247-265); "by construction" vs "approximation,
% not a substitute" from lines 251-252 and 262-265.

The diagnosis came back positive everywhere we looked. Signer identity is
recoverable from every feature space we test, including pose-based ones:
on ASL Citizen the raw spaces separate signers at fixed sign content by a
distance ratio of $22$--$30\times$, and a linear head recovers the signer
far above chance in all three spaces, with the MediaPipe pose features,
often adopted precisely to escape appearance, the most
signer-separable of all, at $98.7\%$ decoding accuracy.
% "recoverable from every feature space we test, including pose-based
% ones" from the abstract, lines 82-84; "$22$--$30\times$" and "a linear
% head recovers the signer far above both chance ... in all three spaces"
% from lines 1355-1359; "98.7%" from line 1348/1361; "most signer-separable"
% from lines 1359-1360.

The two interventions returned two different answers. In the
low-diversity NGT744 setting (four training signers, one studio
environment at train time), synthetic signers more than double
cross-signer retrieval (R@1 $0.23 \to 0.65$ test$\to$train; $0.25 \to
0.70$ train$\to$test, 3 seeds), a gain that a real-multiview control at
the same view count does not reproduce: appearance and identity
diversity, not extra viewpoints, is what drives it. At matched view
count, content, and viewpoint, the generative edits retain roughly
$65$--$80\%$ of the avatar gain, and as real signer diversity is added
at a fixed data budget, the full-strength synthetic gain shrinks,
consistent with synthetic appearance variety substituting for real
diversity that is absent.
% NGT headline numbers and "more than double" from the abstract, lines
% 89-91; multiview control from lines 1525-1534; "65--80%" from lines
% 1518-1520; the shrinking full-Unreal gain (+0.145 -> +0.081 test->train)
% from lines 1611-1614 and fig:ngt_dropoff (lines 1698-1707); "substitutes
% for absent real diversity" phrasing from the fig caption, lines
% 1704-1706. NOTE: the matched-K panel's gain does NOT shrink (lines
% 1615-1621) --- I kept the claim scoped to "full-strength" to stay
% accurate; do not generalise it.

On the large, signer-diverse ASL Citizen benchmark the same idea leaves
recognition essentially unchanged across every objective family we try
(augmentation-as-data, cosine consistency, supervised contrast, and
adversarial invariance), with one modest exception: adding the
counterfactuals to a frozen-feature contrastive head is worth $+1.1$
R@1, while backbone fine-tuning gains nothing over plain cross-entropy.
% "essentially unchanged" + family list from the abstract, lines 93-95;
% "+1.1 R@1 on a frozen-feature head, nothing at the backbone" from line
% 2423; "No appearance or signer-invariance objective beats plain
% cross-entropy at matched depth" from lines 1778-1780.
The embedding-space analyses explain what the retrieval numbers alone
cannot. The counterfactuals themselves are faithful: local edits move an
embedding about half as far as a genuine signer change, the signer swap
exactly as far as one, and although the augments are near-perfectly
linearly separable from real video, that entire linear signature
collapses to a single shared offset per edit type, a direction a
classifier can trivially factor out.
% n.disp 0.44--0.53 local / 1.04--1.05 swap paraphrased from lines
% 2286-2290; "near-perfectly linearly separable" (AUC 0.96--1.00) from
% lines 2305-2310; de-shift collapse to chance (0.48--0.51) and "a
% direction a classifier can trivially factor out" from lines 2311-2323.
And the invariance objectives are not inert: fine-tuning measurably
erodes the linearly decodable component of signer identity,
consistency training most of all, yet that erosion does not translate
into better retrieval, so removing decodable signer information and
improving retrieval accuracy are dissociable outcomes here.
% Decoding erosion (75.3% frozen -> 64.2% CE -> 56.3% consistency) and
% "does not translate into better retrieval ... dissociable outcomes"
% from lines 2092-2105.

The two regimes together give the thesis its answer. Appearance
invariance is not a property a model has or lacks in the abstract; it is
relative to the diversity already present in training. Where real signer
diversity is scarce, as in our four-signer NGT744 corpus and the
few-signer studio corpora of prior positive results, synthetic
appearance counterfactuals earn their cost, and motion-locked
re-renders earn the most; where hundreds of pretraining signers and diverse
recording conditions have already bought the invariance, the same
counterfactuals are largely redundant.
% "not a property a model has or lacks in the abstract; it is relative to
% the diversity already present in training" from lines 2426-2428; the
% prior-positive-results clause refers to the SDDA/PHOENIX comparison
% (nine signers total, five in training), lines 2462-2502; "earn their
% cost ... largely redundant" from the abstract, lines 97-99.
The most direct continuations follow the same axis: parametric avatars
that turn appearance into a sampled variable rather than an authored
asset, a dose-response sweep over the number of synthetic signers, and a
from-scratch control that would test directly whether pretraining
diversity is what consumed the gain in the high-diversity regime.
% Outlook items are the existing future-work themes only: SMPL-X
% parametric avatars (lines 2595-2606), scaling laws / dose-response
% (lines 2608-2611), from-scratch control (lines 2677-2686).
\todome{DRAFT END: (i) no new numbers or claims --- every figure is quoted
from the abstract/results/discussion with line refs in the comments;
(ii) para 3 deliberately scopes the shrinking-gain claim to the
full-strength arm because the matched-$K$ gain stays flat (lines
1615-1621); (iii) para 5's "prior positive results" clause leans on the
SDDA regime comparison in the Discussion --- if you cut or reframe that
section, soften the clause; (iv) I did not cite the pose-diagnosis
implications for the RGB-vs-pose debate here since the Discussion section
for it (line 2533) is still an unresolved todo; (v) the NGT numbers
include the untrimmed-recording noise floor per the Limitations (lines
2542-2549) --- I left that caveat out of the conclusion as the chapter
already owns it, but flag if you want one hedging clause.}
